\documentclass[10pt,a4paper]{amsart}
\usepackage{amsmath,amssymb,amsthm,graphicx}
% Retained preamble for the Stacks Project chapter snapshot in this directory.
% The chapter's own preamble.tex is a copy of the shared Stacks Project preamble; the block
% below is sliced MECHANICALLY, line for line, from it. Nothing was added, removed or
% redefined; the wrapper class, the xr-hyper/hyperref wiring and the externaldocument list
% are omitted because the assembler supplies the document class.
%
% lines taken from preamble.tex: 3, 5, 9, 11, 12, 156, 157, 158, ... (51 lines total)
%
% package coverage restored from the chapter's own preamble (2026-10-13): xy, xy-2cell

\usepackage[all]{xy}
\xyoption{2cell}
\UseAllTwocells

\newenvironment{reference}{\comment}{\endcomment}
\newenvironment{slogan}{\comment}{\endcomment}
\newenvironment{history}{\comment}{\endcomment}
\theoremstyle{plain}
\newtheorem{theorem}[subsection]{Theorem}
\newtheorem{proposition}[subsection]{Proposition}
\newtheorem{lemma}[subsection]{Lemma}
\theoremstyle{definition}
\newtheorem{definition}[subsection]{Definition}
\newtheorem{example}[subsection]{Example}
\newtheorem{exercise}[subsection]{Exercise}
\newtheorem{situation}[subsection]{Situation}
\theoremstyle{remark}
\newtheorem{remark}[subsection]{Remark}
\newtheorem{remarks}[subsection]{Remarks}
\numberwithin{equation}{subsection}
\def\lim{\mathop{\mathrm{lim}}\nolimits}
\def\colim{\mathop{\mathrm{colim}}\nolimits}
\def\Spec{\mathop{\mathrm{Spec}}}
\def\Hom{\mathop{\mathrm{Hom}}\nolimits}
\def\Ext{\mathop{\mathrm{Ext}}\nolimits}
\def\SheafHom{\mathop{\mathcal{H}\!\mathit{om}}\nolimits}
\def\SheafExt{\mathop{\mathcal{E}\!\mathit{xt}}\nolimits}
\def\Sch{\mathit{Sch}}
\def\Mor{\mathop{\mathrm{Mor}}\nolimits}
\def\Ob{\mathop{\mathrm{Ob}}\nolimits}
\def\Sh{\mathop{\mathit{Sh}}\nolimits}
\def\NL{\mathop{N\!L}\nolimits}
\def\CH{\mathop{\mathrm{CH}}\nolimits}
\def\proetale{{pro\text{-}\acute{e}tale}}
\def\etale{{\acute{e}tale}}
\def\QCoh{\mathit{QCoh}}
\def\Ker{\mathop{\mathrm{Ker}}}
\def\Im{\mathop{\mathrm{Im}}}
\def\Coker{\mathop{\mathrm{Coker}}}
\def\Coim{\mathop{\mathrm{Coim}}}
\DeclareMathSymbol{\boxtimes}{\mathbin}{AMSa}{"02}
\def\QCohstack{\mathcal{QC}\!\mathit{oh}}
\def\Cohstack{\mathcal{C}\!\mathit{oh}}
\def\Spacesstack{\mathcal{S}\!\mathit{paces}}
\def\Quotfunctor{\mathrm{Quot}}
\def\Hilbfunctor{\mathrm{Hilb}}
\def\Curvesstack{\mathcal{C}\!\mathit{urves}}
\def\Polarizedstack{\mathcal{P}\!\mathit{olarized}}
\def\Complexesstack{\mathcal{C}\!\mathit{omplexes}}
\def\Pic{\mathop{\mathrm{Pic}}\nolimits}
\def\Picardstack{\mathcal{P}\!\mathit{ic}}
\def\Picardfunctor{\mathrm{Pic}}
\def\Deformationcategory{\mathcal{D}\!\mathit{ef}}
\usepackage{cleveref}
\crefname{originallemma}{original lemma}{original lemmas}
\Crefname{originallemma}{Original lemma}{Original lemmas}
\makeatletter
\newlabel{coherent-lemma-quasi-coherent-affine-cohomology-zero}{{01XB}{source}{}{sourcepointer-1}{}}
\newlabel{coherent-lemma-quasi-coherent-affine-cohomology-zero@cref}{{[originallemma][0][]01XB}{[1][0][]source}{}{}{}}
\makeatother
\crefname{originallemma}{original lemma}{original lemmas}
\Crefname{originallemma}{Original lemma}{Original lemmas}
\makeatletter
\newlabel{descent-lemma-quasi-coherent-and-flat-base-change}{{0756}{source}{}{sourcepointer-2}{}}
\newlabel{descent-lemma-quasi-coherent-and-flat-base-change@cref}{{[originallemma][0][]0756}{[1][0][]source}{}{}{}}
\makeatother
\crefname{originallemma}{original lemma}{original lemmas}
\Crefname{originallemma}{Original lemma}{Original lemmas}
\makeatletter
\newlabel{equation-comparison-modules}{{equation-comparison-modules}{source}{}{sourcepointer-3}{}}
\newlabel{equation-comparison-modules@cref}{{[originallemma][0][]equation-comparison-modules}{[1][0][]source}{}{}{}}
\makeatother
\crefname{originallemma}{original lemma}{original lemmas}
\Crefname{originallemma}{Original lemma}{Original lemmas}
\makeatletter
\newlabel{lemma-cohomology-restriction}{{lemma-cohomology-restriction}{source}{}{sourcepointer-4}{}}
\newlabel{lemma-cohomology-restriction@cref}{{[originallemma][0][]lemma-cohomology-restriction}{[1][0][]source}{}{}{}}
\makeatother
\crefname{originallemma}{original lemma}{original lemmas}
\Crefname{originallemma}{Original lemma}{Original lemmas}
\makeatletter
\newlabel{lemma-comparison}{{lemma-comparison}{source}{}{sourcepointer-5}{}}
\newlabel{lemma-comparison@cref}{{[originallemma][0][]lemma-comparison}{[1][0][]source}{}{}{}}
\makeatother
\crefname{originallemma}{original lemma}{original lemmas}
\Crefname{originallemma}{Original lemma}{Original lemmas}
\makeatletter
\newlabel{lemma-lqc-flat-base-change-fppf-sheaf}{{lemma-lqc-flat-base-change-fppf-sheaf}{source}{}{sourcepointer-6}{}}
\newlabel{lemma-lqc-flat-base-change-fppf-sheaf@cref}{{[originallemma][0][]lemma-lqc-flat-base-change-fppf-sheaf}{[1][0][]source}{}{}{}}
\makeatother
\crefname{originallemma}{original lemma}{original lemmas}
\Crefname{originallemma}{Original lemma}{Original lemmas}
\makeatletter
\newlabel{lemma-pullback-lqc}{{lemma-pullback-lqc}{source}{}{sourcepointer-7}{}}
\newlabel{lemma-pullback-lqc@cref}{{[originallemma][0][]lemma-pullback-lqc}{[1][0][]source}{}{}{}}
\makeatother
\crefname{originallemma}{original lemma}{original lemmas}
\Crefname{originallemma}{Original lemma}{Original lemmas}
\makeatletter
\newlabel{lemma-pushforward-restriction}{{lemma-pushforward-restriction}{source}{}{sourcepointer-8}{}}
\newlabel{lemma-pushforward-restriction@cref}{{[originallemma][0][]lemma-pushforward-restriction}{[1][0][]source}{}{}{}}
\makeatother
\crefname{originallemma}{original lemma}{original lemmas}
\Crefname{originallemma}{Original lemma}{Original lemmas}
\makeatletter
\newlabel{sites-cohomology-lemma-apply-Leray}{{0733}{source}{}{sourcepointer-9}{}}
\newlabel{sites-cohomology-lemma-apply-Leray@cref}{{[originallemma][0][]0733}{[1][0][]source}{}{}{}}
\makeatother
\crefname{originallemma}{original lemma}{original lemmas}
\Crefname{originallemma}{Original lemma}{Original lemmas}
\makeatletter
\newlabel{sites-cohomology-lemma-higher-direct-images}{{072W}{source}{}{sourcepointer-10}{}}
\newlabel{sites-cohomology-lemma-higher-direct-images@cref}{{[originallemma][0][]072W}{[1][0][]source}{}{}{}}
\makeatother
\title{Stacks, Sheaves on Algebraic Stacks: lemma-compare-fppf-etale}
\author{The Stacks Project Authors}
\date{stacks-project:stacks-sheaves.tex:local-snapshot-20261012; GNU FDL 1.2 or later}
\begin{document}
\maketitle
\noindent\textit{Editorial note.} Statement span is lines 4350-4380 of sources/stacks-algebraic-stacks-sheaves/pinned-original-source.tex; proof span is lines 4382-4429, the complete author proof (46 body lines).\par\medskip
\noindent\textit{Editorial note.} Navigation stubs: 10 cited-but-not-included labels are supplied as non-load-bearing pointers with their pinned tags where the snapshot tag index has them.\par\medskip
\noindent\textit{Editorial note.} extra\_packages is empty on purpose: the retained preamble already loads xy with options and a second declaration clashes in the assembler wrapper.\par\medskip
\noindent\textit{Editorial note.} Text is retained verbatim, including the author's own typos and placeholders.\par\medskip
\noindent\textit{Editorial note.} Queue index 11 of queue\_algebraic-stacks-sheaves.json (pb=2048, stmt 4350-4380, proof 4382-4429).\par\medskip
\noindent\textit{Editorial note.} Cross-chapter same-label warning (NOT an occupancy conflict, Chief ruling 2026-10-13 evening): label 'lemma-compare-fppf-etale' also exists in stacks-etale-cohomology — a different proposition in a different chapter file.\par\medskip
\noindent\textit{Allowed standard prerequisites (editorial).}\begin{itemize}
\item Explicit prerequisite boundary. The cited inputs are prerequisites; the argument in the retained span is not replaced by any cited result.
\end{itemize}\medskip
\ifdefined\hypertarget\hypertarget{sourcepointer-1}{}\fi
\noindent\textit{Source reference.} Original lemma 01XB: sources/stacks-algebraic-stacks-sheaves/pinned-original-source.tex, source lines 4429--4429. Cited by the retained statement or proof but not part of the collected spans; supplied only as a navigation pointer. The argument does not depend on its content.\par\medskip
\ifdefined\hypertarget\hypertarget{sourcepointer-2}{}\fi
\noindent\textit{Source reference.} Original lemma 0756: sources/stacks-algebraic-stacks-sheaves/pinned-original-source.tex, source lines 4429--4429. Cited by the retained statement or proof but not part of the collected spans; supplied only as a navigation pointer. The argument does not depend on its content.\par\medskip
\ifdefined\hypertarget\hypertarget{sourcepointer-3}{}\fi
\noindent\textit{Source reference.} Original lemma equation-comparison-modules: sources/stacks-algebraic-stacks-sheaves/pinned-original-source.tex, source lines 1382--1386. Cited by the retained statement or proof but not part of the collected spans; supplied only as a navigation pointer. The argument does not depend on its content.\par\medskip
\ifdefined\hypertarget\hypertarget{sourcepointer-4}{}\fi
\noindent\textit{Source reference.} Original lemma lemma-cohomology-restriction: sources/stacks-algebraic-stacks-sheaves/pinned-original-source.tex, source lines 2788--2803. Cited by the retained statement or proof but not part of the collected spans; supplied only as a navigation pointer. The argument does not depend on its content.\par\medskip
\ifdefined\hypertarget\hypertarget{sourcepointer-5}{}\fi
\noindent\textit{Source reference.} Original lemma lemma-comparison: sources/stacks-algebraic-stacks-sheaves/pinned-original-source.tex, source lines 1303--1334. Cited by the retained statement or proof but not part of the collected spans; supplied only as a navigation pointer. The argument does not depend on its content.\par\medskip
\ifdefined\hypertarget\hypertarget{sourcepointer-6}{}\fi
\noindent\textit{Source reference.} Original lemma lemma-lqc-flat-base-change-fppf-sheaf: sources/stacks-algebraic-stacks-sheaves/pinned-original-source.tex, source lines 4321--4333. Cited by the retained statement or proof but not part of the collected spans; supplied only as a navigation pointer. The argument does not depend on its content.\par\medskip
\ifdefined\hypertarget\hypertarget{sourcepointer-7}{}\fi
\noindent\textit{Source reference.} Original lemma lemma-pullback-lqc: sources/stacks-algebraic-stacks-sheaves/pinned-original-source.tex, source lines 1959--1967. Cited by the retained statement or proof but not part of the collected spans; supplied only as a navigation pointer. The argument does not depend on its content.\par\medskip
\ifdefined\hypertarget\hypertarget{sourcepointer-8}{}\fi
\noindent\textit{Source reference.} Original lemma lemma-pushforward-restriction: sources/stacks-algebraic-stacks-sheaves/pinned-original-source.tex, source lines 4092--4110. Cited by the retained statement or proof but not part of the collected spans; supplied only as a navigation pointer. The argument does not depend on its content.\par\medskip
\ifdefined\hypertarget\hypertarget{sourcepointer-9}{}\fi
\noindent\textit{Source reference.} Original lemma 0733: sources/stacks-algebraic-stacks-sheaves/pinned-original-source.tex, source lines 4429--4429. Cited by the retained statement or proof but not part of the collected spans; supplied only as a navigation pointer. The argument does not depend on its content.\par\medskip
\ifdefined\hypertarget\hypertarget{sourcepointer-10}{}\fi
\noindent\textit{Source reference.} Original lemma 072W: sources/stacks-algebraic-stacks-sheaves/pinned-original-source.tex, source lines 4429--4429. Cited by the retained statement or proof but not part of the collected spans; supplied only as a navigation pointer. The argument does not depend on its content.\par\medskip
\begin{lemma}
\label{lemma-compare-fppf-etale}
Let $S$ be a scheme. Let $\mathcal{X}$ be an algebraic stack over $S$.
Let $\mathcal{F}$ be a presheaf $\mathcal{O}_\mathcal{X}$-module such that
\begin{enumerate}
\item[(a)] $\mathcal{F}$ is locally quasi-coherent, and
\item[(b)] for any morphism $\varphi : x \to y$ of $\mathcal{X}$ which lies
over a morphism of schemes $f : U \to V$ which is flat and
locally of finite presentation, the comparison map
$c_\varphi : f_{small}^*\mathcal{F}|_{V_\etale} \to
\mathcal{F}|_{U_\etale}$ of
(\ref{equation-comparison-modules}) is an isomorphism.
\end{enumerate}
Then $\mathcal{F}$ is an $\mathcal{O}_\mathcal{X}$-module and
we have the following
\begin{enumerate}
\item If $\epsilon : \mathcal{X}_{fppf} \to \mathcal{X}_\etale$
is the comparison morphism, then
$R\epsilon_*\mathcal{F} = \epsilon_*\mathcal{F}$.
\item The cohomology groups $H^p_{fppf}(\mathcal{X}, \mathcal{F})$ are equal
to the cohomology groups computed in the \'etale topology on $\mathcal{X}$.
Similarly for the cohomology groups $H^p_{fppf}(x, \mathcal{F})$ and the
derived versions $R\Gamma(\mathcal{X}, \mathcal{F})$ and
$R\Gamma(x, \mathcal{F})$.
\item If $f : \mathcal{X} \to \mathcal{Y}$ is a $1$-morphism of
categories fibred in groupoids over $(\Sch/S)_{fppf}$ then
$R^if_*\mathcal{F}$ is equal to the fppf-sheafification of the
higher direct image computed in the \'etale cohomology.
Similarly for derived pullback.
\end{enumerate}
\end{lemma}

\begin{proof}
The assertion that $\mathcal{F}$ is an $\mathcal{O}_\mathcal{X}$-module
follows from
Lemma \ref{lemma-lqc-flat-base-change-fppf-sheaf}.
Note that $\epsilon$ is a morphism of sites given by the identity
functor on $\mathcal{X}$. The sheaf $R^p\epsilon_*\mathcal{F}$ is therefore
the sheaf associated to the presheaf
$x \mapsto H^p_{fppf}(x, \mathcal{F})$, see
Cohomology on Sites, Lemma \ref{sites-cohomology-lemma-higher-direct-images}.
To prove (1) it suffices to show that
$H^p_{fppf}(x, \mathcal{F}) = 0$ for $p > 0$
whenever $x$ lies over an affine scheme $U$. By
Lemma \ref{lemma-cohomology-restriction}
we have
$H^p_{fppf}(x, \mathcal{F}) = H^p((\Sch/U)_{fppf}, x^{-1}\mathcal{F})$.
Combining
Descent, Lemma \ref{descent-lemma-quasi-coherent-and-flat-base-change}
with Cohomology of Schemes, Lemma
\ref{coherent-lemma-quasi-coherent-affine-cohomology-zero}
we see that these cohomology groups are zero.

\medskip\noindent
We have seen above that $\epsilon_*\mathcal{F}$ and $\mathcal{F}$ are the
sheaves on $\mathcal{X}_\etale$ and $\mathcal{X}_{fppf}$
corresponding to the same presheaf on $\mathcal{X}$ (and this is true more
generally for any sheaf in the fppf topology on $\mathcal{X}$).
We often abusively identify $\mathcal{F}$ and $\epsilon_*\mathcal{F}$
and this is the sense in which parts (2) and (3) of the lemma should be
understood. Thus part (2) follows formally from (1) and the Leray spectral
sequence, see
Cohomology on Sites, Lemma \ref{sites-cohomology-lemma-apply-Leray}.

\medskip\noindent
Finally we prove (3). The sheaf $R^if_*\mathcal{F}$
(resp.\ $Rf_{\etale, *}\mathcal{F}$)
is the sheaf associated to the presheaf
$$
y \longmapsto
H^i_\tau\Big((\Sch/V)_{fppf} \times_{y, \mathcal{Y}} \mathcal{X},
\ \text{pr}^{-1}\mathcal{F}\Big)
$$
where $\tau$ is $fppf$ (resp.\ $\etale$), see
Lemma \ref{lemma-pushforward-restriction}.
Note that $\text{pr}^{-1}\mathcal{F}$ satisfies properties (a) and (b)
also (by Lemmas \ref{lemma-pullback-lqc} and \ref{lemma-comparison}),
hence these two presheaves are equal by (2).
This immediately implies (3).
\end{proof}

\end{document}
